\documentclass[class=book, crop=false]{standalone}
% ==============================================================================
% SETUP.TEX - CONFIGURACIÓN GLOBAL DEL PROYECTO
% ==============================================================================
% Este archivo centraliza todos los paquetes y estilos.
% Debe incluirse en el preámbulo del libro principal y de los archivos standalone.
% \PassOptionsToPackage{gray}{xcolor}
% --- 1. CODIFICACIÓN E IDIOMA ---
% fix-cm habilita las fuentes Computer Modern en tamaños arbitrarios (por debajo
% de 5 pt, entre otros). Debe cargarse antes que fontenc.
\usepackage{fix-cm}
\usepackage[utf8]{inputenc}
\usepackage[T1]{fontenc}
\usepackage[spanish, es-tabla]{babel}  
\usepackage{csquotes} % Recomendado para babel y citas

\usepackage{import}
% mode=image: Usa el PDF pre-compilado, nunca intenta recompilar.
% Debes compilar cada figura manualmente antes de compilar el libro.
\usepackage[mode=image, subpreambles=false]{standalone}

% --- 2. MATEMÁTICAS Y CIENCIAS ---
\usepackage{amsmath, amssymb, amsfonts, mathtools}
\usepackage{enumitem}

% LaTeX trae \DeclareMathSizes{5}{5}{5}{5}, así que a 5 pt (\tiny) los subíndices
% salen del mismo tamaño que la base: en las etiquetas de figuras el M de
% $\omega_M$ queda desproporcionado. Se restablece la proporción habitual.
\DeclareMathSizes{5}{5}{3.7}{3}

% --- 3. DISEÑO Y GEOMETRÍA --- 
% Unificamos los márgenes para todo el documento
\usepackage[left=2.5cm,right=2.5cm,top=3cm,bottom=3cm]{geometry}
\makeatletter
\ifstandalone % Evita que geometry dañe el recorte de las figuras bajándolas 3cm
  \@ifclasswith{standalone}{crop=false}{
    % Es un capítulo/sección: Dejamos que geometry actúe.
  }{
    % Es una figura: Neutralizamos geometry para que no las recorte mal.
    \geometry{pass}
  }
\fi
\makeatother
\usepackage{parskip} % Espaciado entre párrafos sin sangría
\usepackage{float}   % Para usar [H] en figura

% Ajustes de flotantes para evitar espacios verticales exagerados
\renewcommand{\topfraction}{0.95}      % Permite que las figuras ocupen hasta el 95% superior
\renewcommand{\bottomfraction}{0.95}   % Permite que las figuras ocupen hasta el 95% inferior
\renewcommand{\textfraction}{0.05}     % Permite hojas con tan solo un 5% de texto
\renewcommand{\floatpagefraction}{0.8} % Requiere que una página de solo figuras esté 80% llena
\setcounter{topnumber}{4}
\setcounter{bottomnumber}{4}
\setcounter{totalnumber}{10}

% --- 4. GRÁFICOS Y TABLAS ---
\usepackage{graphicx}
\usepackage{booktabs} % Tablas profesionales
\usepackage{caption}  % Control de captions y \captionof
\usepackage{tikz}
\usetikzlibrary{arrows.meta, calc, angles, quotes, babel, backgrounds}
\usepackage{pgfplots}
\usepgfplotslibrary{groupplots}
\usepgfplotslibrary{external}
\usepgfplotslibrary{patchplots}
\pgfplotsset{compat=1.18}
\usepackage[american, straightvoltages]{circuitikz}

% --- 5. COLORES Y CAJAS ---

\usepackage[table]{xcolor}
\usepackage{tcolorbox}

% Definición de colores corporativos/estilo
\definecolor{utnblue}{RGB}{0, 51, 102}
\definecolor{codegreen}{rgb}{0,0.6,0}
\definecolor{codegray}{rgb}{0.5,0.5,0.5}
\definecolor{codepurple}{rgb}{0.58,0,0.82}
\definecolor{backcolour}{rgb}{0.96,0.96,0.96}

% --- 6. ENTORNOS PERSONALIZADOS (CAJAS) ---

% Caja "Resultado" (Azul Corporativo) — Definiciones, fórmulas clave, resultados importantes.
% Uso: \begin{resultado}[Fórmula de Euler] ... \end{resultado}
\newtcolorbox{resultado}[1][]{
    colback=utnblue!5!white,
    colframe=utnblue,
    title=\textbf{#1},
    fonttitle=\bfseries,
    sharp corners,
    boxrule=0.5mm
}

% Caja "Nota" (Gris) — Advertencias, aplicaciones, contexto secundario.
% Uso: \begin{nota}[Cuidado con la calculadora] ... \end{nota}
% Uso: \begin{nota} ... \end{nota}  → título por defecto "Nota"
\newtcolorbox{nota}[1][Nota]{
    colback=codegray!5!white,
    colframe=codegray,
    title=\textbf{#1},
    fonttitle=\bfseries,
    sharp corners,
    boxrule=0.5mm
}

% Caja inline para resaltar un valor y enlazarlo con su otra aparición en una tabla
% o en una cuenta: mismo color, mismo valor.
% Uso: \marcavalor{$4$}  o  \marcavalor[codegreen]{$4$}
\newtcbox{\marcavalor}[1][utnblue]{
    on line,
    boxsep=1pt, left=2pt, right=2pt, top=1pt, bottom=1pt,
    colback=#1!10!white,
    colframe=#1,
    boxrule=0.4pt,
    arc=2pt
}

% --- 7. CÓDIGO FUENTE (LISTINGS) ---
\usepackage{listings}

\lstdefinestyle{miestilo}{
    backgroundcolor=\color{backcolour},   
    commentstyle=\color{codegreen},
    keywordstyle=\color{magenta},
    numberstyle=\tiny\color{codegray},
    stringstyle=\color{codepurple},
    basicstyle=\ttfamily\footnotesize,
    breakatwhitespace=false,         
    breaklines=true,                 
    captionpos=b,                    
    keepspaces=true,                 
    numbers=left,                    
    numbersep=5pt,                  
    showspaces=false,                
    showstringspaces=false,
    showtabs=false,                  
    tabsize=2,
    frame=single,                    
    rulecolor=\color{codegray}
}

\lstset{style=miestilo}
% Soporte para tildes y ñ en bloques de código
\lstset{literate={ñ}{{\~n}}1 {á}{{\'a}}1 {é}{{\'e}}1 {í}{{\'i}}1 {ó}{{\'o}}1 {ú}{{\'u}}1}

% Operadores matemáticos personalizados

% Vector en sentido discreto (lista finita de coordenadas): negrita + flecha.
% Uso: \vect{v}, \vect{e}_k, \vect{0}.
\newcommand{\vect}[1]{\vec{\mathbf{#1}}}

% Fecha de versión y comando para capítulos standalone
\providecommand{\verdate}{\the\year.\ifnum\month<10 0\fi\the\month.\ifnum\day<10 0\fi\the\day}
\newcommand{\chapterversion}{%
    \vspace{-1.5cm}\begin{flushright}\small\textit{\color{codegray}Versión~\verdate}\end{flushright}\vspace{0.5cm}%
}

% --- 8. HIPERVÍNCULOS (DEBE SER EL ÚLTIMO PAQUETE) ---
\usepackage[colorlinks=true, linkcolor=utnblue, urlcolor=utnblue, citecolor=utnblue]{hyperref}

% --- 9. REFERENCIAS CRUZADAS ENTRE CAPÍTULOS STANDALONE ---
% Cuando se compila un capítulo standalone, xr-hyper lee los .aux de los
% otros capítulos (si existen) para resolver \ref{} a labels externos.
% En la compilación del libro completo este bloque no se activa.
\ifstandalone
  \usepackage{xr-hyper}
  \makeatletter
  \newcommand{\xrchap}[1]{%
    \edef\xrc@self{\detokenize{Capitulo#1}}%
    \edef\xrc@job{\jobname}%
    \ifx\xrc@self\xrc@job\else
      \IfFileExists{../#1capitulo/Capitulo#1.aux}%
        {\externaldocument{../#1capitulo/Capitulo#1}}{}%
    \fi
  }
  \makeatother
  \xrchap{01}\xrchap{02}\xrchap{03}\xrchap{04}
  \xrchap{05}\xrchap{06}\xrchap{07}\xrchap{08}
  \xrchap{09}\xrchap{10}\xrchap{11}\xrchap{12}
  \xrchap{13}
\fi

\begin{document}
\setcounter{chapter}{5}
\section{Serie de Fourier exponencial y armónica}
\label{sec:serie_fourier}

En esta sección construimos la representación de una señal
\textbf{periódica} $x(t)$, de período $T_0$ y frecuencia angular
$\omega_0 = 2\pi/T_0$, como combinación lineal de las exponenciales
armónicas $\{e^{jk\omega_0 t}\}_{k\in\mathbb{Z}}$. La sección anterior
dejó las dos piezas que necesitamos para hacerlo: el producto interno
entre señales y la verificación de que esas exponenciales son
ortogonales sobre cualquier intervalo de longitud $T_0$. Con esos dos
ingredientes, la fórmula de descomposición sobre base ortogonal de
$\mathbb{R}^n$ ---ecuación~\eqref{eq:coefs_ortogonal_Rn}--- se traslada
literalmente al espacio de señales: dada una base ortogonal
$\{\phi_k\}$ y una señal $x(t)$, los coeficientes de la combinación
$x(t) = \sum_k c_k\,\phi_k(t)$ son
\begin{equation}
    c_k \;=\; \frac{\langle x,\phi_k\rangle}{\|\phi_k\|^2}.
    \label{eq:coefs_generalizados}
\end{equation}
Lo único que cambia entre el caso finito y el de señales es la forma
del producto interno: en $\mathbb{R}^n$ era una suma sobre coordenadas;
en señales es la integral $\int x(t)\,\overline{\phi_k(t)}\,dt$
definida en~\eqref{eq:pi_senales}. La estructura es idéntica. Lo que
sigue del capítulo aplica esta fórmula al caso particular en que los
$\phi_k$ son las exponenciales armónicas. Ahí los coeficientes dejan
de ser etiquetas abstractas y empiezan a hablar de \emph{frecuencia}.

% ==============================================================================
% 5.2.1 Serie exponencial compleja
% ==============================================================================
\subsection{Serie exponencial compleja}
\label{subsec:serie_exp_compleja}

Aplicamos la receta~\eqref{eq:coefs_generalizados} al sistema de
exponenciales armónicas $\phi_k(t) = e^{jk\omega_0 t}$ sobre un
período. La norma cuadrática de cada elemento de la base es inmediata:
con $n = m = k$ en la ortogonalidad exponencial~\eqref{eq:ort_exp},
\begin{equation*}
    \|\phi_k\|^2
    \;=\; \int_{T_0} e^{jk\omega_0 t}\,\overline{e^{jk\omega_0 t}}\,dt
    \;=\; \int_{T_0} 1\,dt
    \;=\; T_0,
\end{equation*}
independientemente de $k$. Y el producto interno de $x$ con $\phi_k$,
recordando que $\overline{e^{jk\omega_0 t}} = e^{-jk\omega_0 t}$, es
\begin{equation*}
    \langle x,\phi_k\rangle
    \;=\; \int_{T_0} x(t)\,e^{-jk\omega_0 t}\,dt.
\end{equation*}
Reemplazando en~\eqref{eq:coefs_generalizados}, los coeficientes
quedan
\begin{equation}
    c_k \;=\; \frac{1}{T_0}\int_{T_0} x(t)\,e^{-jk\omega_0 t}\,dt,
    \qquad k \in \mathbb{Z},
    \label{eq:ck_analisis}
\end{equation}
y la combinación lineal correspondiente,
\begin{equation}
    x(t) \;=\; \sum_{k=-\infty}^{\infty} c_k\,e^{jk\omega_0 t},
    \label{eq:sf_sintesis}
\end{equation}
es la \textbf{serie de Fourier exponencial} de la señal periódica
$x(t)$. La igualdad en~\eqref{eq:sf_sintesis} requiere precisar en
qué sentido se entiende ---convergencia puntual, en norma cuadrática,
comportamiento en las discontinuidades---, y eso es el tema de la
sección siguiente; en lo que resta de la actual la tomamos como
asumida. La fórmula~\eqref{eq:ck_analisis} se llama de \emph{análisis}
---calcula cada $c_k$ a partir de la señal--- y
la~\eqref{eq:sf_sintesis} de \emph{síntesis} ---reconstruye la señal a
partir de los coeficientes---.

Hasta aquí la receta. Lo que vuelve esto interpretable, y lo que
conecta la serie de Fourier con lo construido en el
Capítulo~\ref{cap02}, es leer cada término $c_k\,e^{jk\omega_0 t}$
geométricamente. La exponencial $e^{jk\omega_0 t}$ es el
\textbf{vector giratorio} de~\eqref{eq:vector_giratorio} con
$\omega = k\omega_0$: para cada $t$, un punto de módulo $1$ en el
plano complejo; como función de $t$, recorre la circunferencia
unitaria a velocidad angular $k\omega_0$. Multiplicarlo por el
coeficiente complejo $c_k = |c_k|\,e^{j\angle c_k}$ escala su
módulo a $|c_k|$ y desplaza su fase inicial en $\angle c_k$, sin
alterar la velocidad de giro. Cada término $c_k\,e^{jk\omega_0 t}$
de la síntesis es entonces una versión escalada y desfasada del
vector giratorio $e^{jk\omega_0 t}$.

Conviene leer la suma~\eqref{eq:sf_sintesis} término a término.
Desplegándola alrededor del cero,
\begin{equation*}
    x(t) \;=\; \cdots
            + c_{-2}\,e^{-j 2\omega_0 t}
            + c_{-1}\,e^{-j\omega_0 t}
            + c_{0}
            + c_{1}\,e^{j\omega_0 t}
            + c_{2}\,e^{j 2\omega_0 t}
            + \cdots,
\end{equation*}
a cada sumando $c_k\,e^{jk\omega_0 t}$ lo llamamos
\textbf{armónico} de la señal y lo nombramos por su orden $k$:
el armónico de $k=1$, de frecuencia $\omega_0$ ---la del propio
período $T_0$---, es el \textbf{armónico fundamental} (o
simplemente la \emph{fundamental}); el de $k=2$, el
\textbf{segundo armónico}; el de $k=3$, el \textbf{tercer
armónico}; y así sucesivamente. La frecuencia $\omega_0$ se
llama también \textbf{frecuencia fundamental} de la señal, y los
múltiplos enteros $k\omega_0$ son las frecuencias de los
armónicos sucesivos. La síntesis reconstruye $x(t)$ como
superposición de todos sus armónicos, uno por cada
$k\in\mathbb{Z}$.

El armónico $k=0$ es la única excepción al giro. Su exponencial
vale $e^{j\cdot 0\cdot \omega_0 t} = 1$ ---no rota, no oscila,
no depende del tiempo---, así que el sumando se reduce a la
constante $c_0$. Es el único armónico constante de toda la
suma; los demás oscilan en el tiempo. Su valor sale directo de
la fórmula de análisis~\eqref{eq:ck_analisis} con $k=0$, donde
el integrando se reduce a $x(t)$:
\begin{equation*}
    c_0 \;=\; \frac{1}{T_0}\int_{T_0} x(t)\,dt.
\end{equation*}
El segundo miembro es exactamente el \textbf{valor medio} de
$x(t)$ sobre un período ---introducido en el
Capítulo~\ref{cap03} y también llamado \emph{componente
continua} o \emph{componente DC}---. La serie de Fourier
descompone entonces a la señal en un valor medio $c_0$,
alrededor del cual no hay oscilación, más la superposición de
los armónicos $c_k\,e^{jk\omega_0 t}$ con $k\neq 0$ que sí
oscilan.

El índice $k$ corre sobre todos los enteros, positivos y
negativos. El vector giratorio $e^{jk\omega_0 t}$ con $k > 0$
rota antihorariamente; con $k < 0$, en sentido horario, según
la convención de signo del Capítulo~\ref{cap02}. Cuando $x(t)$
es real ---el caso usual del libro---, los coeficientes quedan
ligados por la simetría conjugada $c_{-k} = \overline{c_k}$, y
entonces las componentes $c_k\,e^{jk\omega_0 t}$ y
$c_{-k}\,e^{-jk\omega_0 t}$ son vectores giratorios conjugados
---uno antihorario, uno horario, con módulos iguales---, cuya
suma es una sinusoide real, como vimos en el
Capítulo~\ref{cap02}. La formalizaremos como propiedad en la
Subsección~\ref{subsec:prop_conjugacion}.

\begin{resultado}[Serie de Fourier exponencial]
Sea $x(t)$ una señal periódica de período $T_0$ y frecuencia angular
$\omega_0 = 2\pi/T_0$. Sus coeficientes de Fourier son
\begin{equation}
    c_k \;=\; \frac{1}{T_0}\int_{T_0} x(t)\,e^{-jk\omega_0 t}\,dt,
    \qquad k \in \mathbb{Z},
    \label{eq:ck_analisis_box}
\end{equation}
y la síntesis de la señal a partir de ellos es
\begin{equation}
    x(t) \;=\; \sum_{k=-\infty}^{\infty} c_k\,e^{jk\omega_0 t}.
    \label{eq:sf_sintesis_box}
\end{equation}
\end{resultado}

\paragraph{Señales complejas y reales.}
La fórmula de análisis~\eqref{eq:ck_analisis_box} acepta
cualquier señal periódica $x(t)$, sea esta real o compleja; las
afirmaciones de las secciones siguientes que dependan de la
hipótesis ``$x$ real'' las marcaremos explícitamente.

\begin{nota}[Animaciones del capítulo]
Las exponenciales armónicas son objetos en movimiento, y la imagen
estática del libro pierde parte de la geometría. Junto al capítulo
se incluyen animaciones que muestran el giro de los vectores
$c_k\,e^{jk\omega_0 t}$ y la formación de la señal periódica como
superposición de pares conjugados que giran en sentidos opuestos.
\end{nota}

% ==============================================================================
% 5.2.2 Ejemplo: pulso rectangular periódico
% ==============================================================================
\subsection{Ejemplo: pulso rectangular periódico}
\label{subsec:ejemplo_pulso}

Aplicamos la receta a un \textbf{tren periódico de pulsos
rectangulares} con valores numéricos concretos: la señal toma
amplitud $1$ sobre una ventana centrada de ancho $\tau = 1$ y vale
$0$ en el resto del período $T_0 = 4$,
\begin{equation*}
    x(t) \;=\;
    \begin{cases}
        1, & |t| \leq 1/2, \\
        0, & 1/2 < |t| \leq 2,
    \end{cases}
    \qquad
    x(t + 4) \;=\; x(t),
\end{equation*}
con frecuencia angular $\omega_0 = 2\pi/T_0 = \pi/2$. La
figura~\ref{fig:tren_pulsos} grafica dos períodos de la señal. Para
aprovechar la simetría par del pulso integramos sobre el intervalo
simétrico $[-2,\,2]$.

\begin{figure}[H]
    \centering
    \begin{minipage}{\linewidth}
        \centering
        \includestandalone[width=0.75\linewidth]{figures/fig_tren_pulsos/fig_tren_pulsos}
        \caption{Tren periódico de pulsos rectangulares de período
                 $T_0 = 4$ y ancho $\tau = 1$. La amplitud vale $1$
                 sobre la ventana centrada en cada múltiplo de $T_0$
                 y $0$ entre ventanas.}
        \label{fig:tren_pulsos}
    \end{minipage}
\end{figure}

\paragraph{Cálculo de $c_0$.}
Por la lectura de $c_0$ como valor medio de la señal, el
coeficiente sale por inspección: el pulso vale $1$ sobre una
ventana de ancho $1$ y $0$ sobre el resto del período de
longitud $T_0=4$, así que su promedio es la cuarta parte de su
amplitud,
\begin{equation*}
    c_0 \;=\; \frac{1}{4}.
\end{equation*}
Aplicar la fórmula de análisis~\eqref{eq:ck_analisis_box} con
$k=0$ ---el camino general, útil cuando la inspección no es tan
directa--- da el mismo número:
\begin{equation*}
    c_0 \;=\; \frac{1}{T_0}\int_{T_0} x(t)\,dt
        \;=\; \frac{1}{4}\int_{-1/2}^{1/2} 1\,dt
        \;=\; \frac{1}{4}.
\end{equation*}

\paragraph{Cálculo de $c_k$ para $k \neq 0$.}
Como $x(t) = 1$ dentro de la ventana central y $0$ fuera, la
integral de análisis se reduce a la ventana:
\begin{equation*}
    c_k \;=\; \frac{1}{4}\int_{-1/2}^{1/2} e^{-jk\omega_0 t}\,dt.
\end{equation*}
La primitiva de $e^{-jk\omega_0 t}$ es
$e^{-jk\omega_0 t}/(-jk\omega_0)$. Evaluando entre $-1/2$ y $1/2$,
\begin{equation*}
    c_k \;=\; \frac{1}{4}\cdot
        \frac{e^{-jk\omega_0/2} - e^{jk\omega_0/2}}{-jk\omega_0}.
\end{equation*}
Conviene reordenar el numerador en el orden estándar de la
identidad de Euler $e^{j\theta} - e^{-j\theta} = 2j\sin\theta$.
Para eso intercambiamos los dos términos de la resta y absorbemos
el signo negativo en el prefactor:
\begin{equation*}
    c_k \;=\; \frac{1}{4}\cdot
        \frac{e^{jk\omega_0/2} - e^{-jk\omega_0/2}}{jk\omega_0}.
\end{equation*}
Aplicando ahora la identidad de Euler con $\theta = k\omega_0/2$,
\begin{equation*}
    c_k \;=\; \frac{1}{4}\cdot
        \frac{2j\sin(k\omega_0/2)}{jk\omega_0}
    \;=\; \frac{\sin(k\omega_0/2)}{2\,k\omega_0},
\end{equation*}
donde el factor $j$ del numerador canceló contra el del
denominador. Para reconocer el seno cardinal
$\operatorname{sinc}(x) = \sin(x)/x$ ---introducido
en la Subsección~\ref{subsec:sinc}--- conviene un último paso de cosmética:
multiplicar numerador y denominador por $2$ de modo que el
argumento del seno aparezca también en el denominador,
\begin{equation*}
    c_k \;=\; \frac{1}{4}\cdot
        \frac{\sin(k\omega_0/2)}{k\omega_0/2}
        \;=\; \frac{1}{4}\,\operatorname{sinc}\!\Bigl(\frac{k\omega_0}{2}\Bigr).
\end{equation*}
Sustituyendo el valor concreto $\omega_0 = \pi/2$ del ejemplo,
\begin{equation}
    c_k \;=\; \frac{1}{4}\,\operatorname{sinc}\!\Bigl(\frac{k\pi}{4}\Bigr),
    \qquad k \in \mathbb{Z}.
    \label{eq:ck_pulso}
\end{equation}
La fórmula vale también para $k=0$ con la convención
$\operatorname{sinc}(0) = 1$: recupera $c_0 = 1/4$ y deja una sola
expresión para todos los enteros. Se ve ahora por qué convenía
calcular $c_0$ aparte: la primitiva $e^{-jk\omega_0 t}/(-jk\omega_0)$
de la que partimos tiene $k$ en el denominador, y la cuenta entera
no estaba definida en $k=0$. La sinc al final unifica los dos casos,
pero para llegar a ella hubo que pasar por una primitiva que excluía
el cero.

\paragraph{Coeficientes uno a uno.}
Evaluemos~\eqref{eq:ck_pulso} en los primeros enteros, agregando
en cada fila la \textbf{forma polar} ---\emph{módulo} $\angle$
\emph{fase}---, que es la que hace falta para leer cada
coeficiente como un vector en el plano. Como los $c_k$ del pulso
son reales, la fase vale $0$ cuando $c_k > 0$ y $\pi$ cuando
$c_k < 0$ (porque $-|c_k| = |c_k|\,e^{j\pi}$). El armónico
$k=0$ va aparte; el resto lo agrupamos en columnas paralelas,
$c_k$ a la izquierda y $c_{-k}$ a la derecha:
\begin{align*}
    c_0  &\;=\; \tfrac{1}{4}\,\operatorname{sinc}(0)
                \;=\; \tfrac{1}{4}
                \;=\; \tfrac{1}{4}\,\angle\,0, \\[2pt]
    c_1  &\;=\; \tfrac{1}{4}\,\operatorname{sinc}(\pi/4)
                \;\approx\; 0{,}225
                \;\approx\; 0{,}225\,\angle\,0,
    & c_{-1} &\;=\; \tfrac{1}{4}\,\operatorname{sinc}(-\pi/4)
                \;\approx\; 0{,}225
                \;\approx\; 0{,}225\,\angle\,0, \\[2pt]
    c_2  &\;=\; \tfrac{1}{4}\,\operatorname{sinc}(\pi/2)
                \;\approx\; 0{,}159
                \;\approx\; 0{,}159\,\angle\,0,
    & c_{-2} &\;=\; \tfrac{1}{4}\,\operatorname{sinc}(-\pi/2)
                \;\approx\; 0{,}159
                \;\approx\; 0{,}159\,\angle\,0, \\[2pt]
    c_3  &\;=\; \tfrac{1}{4}\,\operatorname{sinc}(3\pi/4)
                \;\approx\; 0{,}075
                \;\approx\; 0{,}075\,\angle\,0,
    & c_{-3} &\;=\; \tfrac{1}{4}\,\operatorname{sinc}(-3\pi/4)
                \;\approx\; 0{,}075
                \;\approx\; 0{,}075\,\angle\,0, \\[2pt]
    c_4  &\;=\; \tfrac{1}{4}\,\operatorname{sinc}(\pi)
                \;=\; 0,
    & c_{-4} &\;=\; \tfrac{1}{4}\,\operatorname{sinc}(-\pi)
                \;=\; 0, \\[2pt]
    c_5  &\;=\; \tfrac{1}{4}\,\operatorname{sinc}(5\pi/4)
                \;\approx\; -0{,}045
                \;\approx\; 0{,}045\,\angle\,\pi,
    & c_{-5} &\;=\; \tfrac{1}{4}\,\operatorname{sinc}(-5\pi/4)
                \;\approx\; -0{,}045
                \;\approx\; 0{,}045\,\angle\,\pi.
\end{align*}
Las dos columnas coinciden número a número: $c_{-k} = c_k$
porque el seno cardinal es par en su argumento.

Dos observaciones se desprenden del listado, y vale la pena
nombrarlas porque vamos a retomar este ejemplo en la
Sección~\ref{sec:convergencia_propiedades}. Primero, los $c_k$ del pulso
resultan \textbf{reales} para todo $k$ ---en polar, las fases solo
toman los valores $0$ y $\pi$, sin valores intermedios---. Segundo,
sus módulos son \textbf{pares} en $k$ ---$|c_{-k}| = |c_k|$--- y la
sucesión cambia de signo en cada cero del seno cardinal: positiva para
$|k| \leq 3$, nula en $|k| = 4$, negativa entre $|k|=5$ y $|k|=7$,
y así sucesivamente. Las dos observaciones no son casualidad: en la
sección siguiente vamos a ver que son consecuencias directas de la
simetría conjugada de las señales reales y de la simetría par del
pulso.

\paragraph{Lectura geométrica de los términos.}
Cada término $c_k\,e^{jk\omega_0 t}$ de la
síntesis~\eqref{eq:sf_sintesis_box} es un vector giratorio de
módulo $|c_k|$ y velocidad angular $k\omega_0$, en la lectura
desarrollada en la Subsección~\ref{subsec:serie_exp_compleja}. Para nuestro
ejemplo, la velocidad angular del armónico $k$ es
$k\omega_0 = k\pi/2$, y los módulos $|c_k|$ son los del listado en
forma polar.

La figura~\ref{fig:vectores_pulso} muestra ocho vectores
---correspondientes a $k = \pm 1, \pm 2, \pm 3, \pm 4$---: los de
$k>0$ giran antihorariamente sobre circunferencias de radio
$|c_k|$, y los de $k<0$ giran horariamente sobre circunferencias
del mismo radio (los módulos son pares en $k$, como acabamos de
ver). La superposición de estos vectores, junto con los infinitos
restantes, reconstruye la señal real $x(t)$.

La fase ya quedó leída en el listado polar. Para $|k| \leq 3$ vale
$0$: el vector arranca apuntando al semieje real positivo. Recién
a partir de $|k|=5$ ---donde el seno cardinal cruza por su primer cero y
vuelve con signo opuesto--- la fase pasa a $\pi$, y los vectores
arrancan apuntando al semieje real negativo (rotación inicial de
media vuelta).

\begin{figure}[H]
    \centering
    \begin{minipage}{\linewidth}
        \centering
        \includestandalone[width=\linewidth]{figures/fig_vectores_pulso/fig_vectores_pulso}
        \caption{Vectores giratorios $c_k\,e^{jk\omega_0 t}$ del
                 tren de pulsos rectangulares con $\tau = 1$,
                 $T_0 = 4$. Cada panel muestra la circunferencia
                 de radio $|c_k|$ y el sentido de giro, con el
                 orden $k$ del armónico marcado arriba. La fila
                 superior corresponde a los $k$ positivos (giro
                 antihorario) y la inferior a los negativos (giro
                 horario). El armónico $k = \pm 4$ tiene módulo
                 nulo porque $\operatorname{sinc}(\pi) = 0$.}
        \label{fig:vectores_pulso}
    \end{minipage}
\end{figure}

\paragraph{Función de síntesis.}
Reuniendo todo, la serie de Fourier del pulso tiene en forma
compacta la expresión
\begin{equation*}
    x(t) \;=\; \sum_{k=-\infty}^{\infty}
        \frac{1}{4}\,\operatorname{sinc}\!\Bigl(\frac{k\pi}{4}\Bigr)
        \,e^{j k \pi t/2},
\end{equation*}
y, sustituyendo los valores numéricos de los primeros armónicos,
se despliega como
\begin{align*}
    x(t) \;\approx\; \cdots
        & + 0{,}075\,e^{-j 3\pi t/2}
        + 0{,}159\,e^{-j\pi t}
        + 0{,}225\,e^{-j\pi t/2}
        + \tfrac{1}{4} \\
        & + 0{,}225\,e^{j\pi t/2}
        + 0{,}159\,e^{j\pi t}
        + 0{,}075\,e^{j 3\pi t/2}
        + \cdots
\end{align*}
La señal queda escrita como un nivel medio $1/4$ más una
superposición simétrica de armónicos en pares conjugados.

\paragraph{Reconstrucción truncada.}
Recuperamos la señal aplicando la síntesis~\eqref{eq:sf_sintesis_box},
pero truncando la suma a un número finito de armónicos:
\begin{equation*}
    \hat{x}_N(t) \;=\; \sum_{k=-N}^{N} c_k\,e^{jk\omega_0 t}.
\end{equation*}
La figura~\ref{fig:pulso_reconstruccion} muestra $\hat{x}_N(t)$ para
$N = 1, 3, 7, 21$ superpuesta al pulso original. Al aumentar $N$ la
aproximación se ajusta cada vez mejor sobre los tramos donde la señal
es constante, pero \emph{persiste un sobrepico junto a las
discontinuidades} que no se atenúa al aumentar $N$: solo se vuelve más
angosto y se desplaza más cerca del salto. Este comportamiento
puntual, conocido como \emph{fenómeno de Gibbs}, queda completamente
explicado en la sección siguiente.

\begin{figure}[H]
    \centering
    \begin{minipage}{\linewidth}
        \centering
        \includestandalone[width=0.7\linewidth]{figures/fig_pulso_reconstruccion/fig_pulso_reconstruccion}
        \caption{Aproximación $\hat{x}_N(t) = \sum_{k=-N}^{N} c_k
                 e^{jk\omega_0 t}$ del tren de pulsos rectangulares
                 ($\tau = 1$, $T_0 = 4$). En cada panel, la señal
                 original en azul y la aproximación truncada en rojo.
                 (a) $N = 1$. (b) $N = 3$. (c) $N = 7$. (d) $N = 21$.
                 Al aumentar $N$ la curva se acerca a la señal en los
                 tramos constantes, pero el sobrepico junto a las
                 discontinuidades no se atenúa.}
        \label{fig:pulso_reconstruccion}
    \end{minipage}
\end{figure}

% ==============================================================================
% 5.2.3 Espectros de líneas
% ==============================================================================
\subsection{Espectros de líneas}
\label{subsec:espectros_lineas}

Los coeficientes $c_k$ son números complejos indexados por un
entero, y como tales se grafican naturalmente como dos
sucesiones discretas: el \textbf{espectro de módulos} $|c_k|$
y el \textbf{espectro de fases} $\angle c_k$. Cada coeficiente
ocupa un punto de la \emph{grilla armónica} ---el eje
horizontal discreto, con un lugar reservado para cada
$k\in\mathbb{Z}$--- y se dibuja como un \emph{stem}
---segmento vertical con un punto en el extremo---. La imagen
resultante, llamada \textbf{espectro de líneas}, es la
representación gráfica más usada de una señal periódica en el
dominio de la frecuencia.

El eje horizontal admite dos convenciones equivalentes para
rotularse: por la frecuencia $k\omega_0$ ---grilla separada
por $\omega_0$, con unidades físicas--- o por el índice entero
$k$ ---grilla unitaria, números más limpios---. Las dos
aparecen en la práctica; lo que importa es elegir una y
mantenerla dentro de un mismo gráfico. En esta subsección
usamos el índice $k$.

El espectro de líneas determina por completo a $x(t)$ vía la
fórmula de síntesis~\eqref{eq:sf_sintesis_box}: dos
representaciones equivalentes de la misma señal, una sobre el
eje del tiempo y otra sobre la grilla armónica.

\paragraph{Espectro del pulso rectangular.}
Con los $c_k$ ya calculados, el espectro de líneas del pulso
sale por simple graficación. La
figura~\ref{fig:espectro_pulso} muestra el módulo $|c_k|$ y la
fase $\angle c_k$ contra el índice $k$.

\begin{figure}[H]
    \centering
    \begin{minipage}{\linewidth}
        \centering
        \includestandalone[width=0.65\linewidth]{figures/fig_espectro_pulso/fig_espectro_pulso}
        \caption{Espectro de líneas del tren de pulsos rectangulares
                 con $\tau = 1$, $T_0 = 4$.
                 (a) Módulo $|c_k|$ con la envolvente sinc en línea
                 punteada; los ceros aparecen en los múltiplos enteros
                 de $T_0/\tau = 4$.
                 (b) Fase $\angle c_k$, que alterna entre $0$ y $\pi$.}
        \label{fig:espectro_pulso}
    \end{minipage}
\end{figure}

Mirando los gráficos hay un hecho que sumar a las observaciones
del listado: la fase $\angle c_k$ es \textbf{impar} en $k$
---los stems de $k$ y $-k$ son opuestos en signo---. Es otra
firma gráfica de las señales periódicas reales, que se
formaliza junto con la paridad del módulo en la
Subsección~\ref{subsec:prop_conjugacion}.

% ==============================================================================
% 5.2.4 Ejemplo: diente de sierra
% ==============================================================================
\subsection{Ejemplo: diente de sierra}
\label{subsec:ejemplo_sierra}

Como segundo ejemplo desarrollaremos el \textbf{diente de
sierra}. La señal crece linealmente desde $0$ hasta $2$ a lo
largo de un período y vuelve a $0$ por un salto en cada
múltiplo entero del período. Trabajamos con el mismo período
$T_0 = 4$ del pulso ---y la misma frecuencia fundamental
$\omega_0 = \pi/2$---, ubicamos el salto en $t = 0$ y
definimos
\begin{equation*}
    x(t) \;=\; \frac{t}{2}, \qquad t \in [0, 4),
    \qquad x(t+4) = x(t).
\end{equation*}
La figura~\ref{fig:diente_sierra_sf} grafica dos períodos.

\begin{figure}[H]
    \centering
    \begin{minipage}{\linewidth}
        \centering
        \includestandalone[width=0.75\linewidth]{figures/fig_diente_sierra/fig_diente_sierra}
        \caption{Diente de sierra de período $T_0 = 4$. La
                 rampa sube de $0$ a $2$ dentro de cada
                 período y vuelve a $0$ por el salto en los
                 múltiplos enteros de $T_0$.}
        \label{fig:diente_sierra_sf}
    \end{minipage}
\end{figure}

A diferencia del pulso, el diente de sierra crece de manera
monótona y termina en un salto. Tiene \textbf{nivel medio} $1$
---el promedio entre el inicio en $0$ y el final en $2$---, y
oscila alrededor de ese valor.

\paragraph{Cálculo de $c_0$.}
$c_0$ es el valor medio de la señal sobre un período. Para la
rampa basta mirar la figura~\ref{fig:diente_sierra_sf}: como sube
linealmente desde $0$ hasta $2$ a lo largo del período, su
promedio es el punto medio entre los dos extremos,
$c_0 = (0+2)/2 = 1$. Dicho de otro modo, la línea horizontal
$y = 1$ cruza la rampa en la mitad del período y deja igual área
arriba y abajo de ella.

La cuenta analítica lo confirma. Tomando $k=0$ en la fórmula de
análisis~\eqref{eq:ck_analisis_box},
\begin{equation*}
    c_0 \;=\; \frac{1}{4}\int_{0}^{4} \frac{t}{2}\,dt
        \;=\; \frac{1}{8}\int_{0}^{4} t\,dt
        \;=\; \frac{1}{8}\left[\frac{t^2}{2}\right]_{0}^{4}
        \;=\; \frac{1}{8}\cdot 8
        \;=\; 1.
\end{equation*}

\paragraph{Cálculo de $c_k$ para $k \neq 0$.}
Aplicando la fórmula de análisis~\eqref{eq:ck_analisis_box}
sobre el período $[0, T_0)$,
\begin{equation*}
    c_k \;=\; \frac{1}{4}\int_{0}^{4} \frac{t}{2}\,
            e^{-jk\omega_0 t}\,dt
        \;=\; \frac{1}{8}\int_{0}^{4} t\,e^{-jk\omega_0 t}\,dt.
\end{equation*}
La integral pide \emph{integración por partes}, con $u = t$ y
$dv = e^{-jk\omega_0 t}\,dt$:
\begin{equation*}
    \int_{0}^{4} t\,e^{-jk\omega_0 t}\,dt
    \;=\; \left[\frac{t\,e^{-jk\omega_0 t}}{-jk\omega_0}\right]_{0}^{4}
    \;-\; \int_{0}^{4} \frac{e^{-jk\omega_0 t}}{-jk\omega_0}\,dt.
\end{equation*}
La segunda integral se anula: el integrando es una exponencial
armónica de orden $k\neq 0$ y se integra sobre un período
completo. Queda solo el término de borde. Como
$e^{-jk\omega_0 \cdot 4} = e^{-2jk\pi} = 1$ para todo entero
$k$,
\begin{equation*}
    \left[\frac{t\,e^{-jk\omega_0 t}}{-jk\omega_0}\right]_{0}^{4}
    \;=\; \frac{4\cdot 1 - 0}{-jk\omega_0}
    \;=\; \frac{4}{-jk\omega_0}
    \;=\; \frac{4j}{-j^2\,k\omega_0}
    \;=\; \frac{4j}{k\omega_0},
\end{equation*}
donde multiplicamos numerador y denominador por $j$ y usamos
$j^2 = -1$. Sustituyendo en $c_k$,
\begin{equation*}
    c_k \;=\; \frac{1}{8}\cdot\frac{4j}{k\omega_0}
        \;=\; \frac{j}{2\,k\omega_0},
\end{equation*}
y reemplazando $\omega_0 = \pi/2$,
\begin{equation}
    c_k \;=\; \frac{j}{\pi k},
        \qquad k \neq 0.
    \label{eq:ck_sierra}
\end{equation}

\paragraph{Coeficientes uno a uno.}
Los $c_k$ del diente de sierra son \textbf{puramente
imaginarios} para $k\neq 0$: la fase vale $+\pi/2$ para los
$k>0$ y $-\pi/2$ para los $k<0$, y el módulo $1/(\pi|k|)$
decrece como $1/|k|$. Listando los primeros enteros en pares
de índice opuesto, ya en forma polar:
\begin{align*}
    c_0 &\;=\; 1, \\[2pt]
    c_1 &\;=\; j/\pi
            \;\approx\; 0{,}318\,\angle\,(\pi/2),
    & c_{-1} &\;=\; -j/\pi
            \;\approx\; 0{,}318\,\angle\,(-\pi/2), \\[2pt]
    c_2 &\;=\; j/(2\pi)
            \;\approx\; 0{,}159\,\angle\,(\pi/2),
    & c_{-2} &\;=\; -j/(2\pi)
            \;\approx\; 0{,}159\,\angle\,(-\pi/2), \\[2pt]
    c_3 &\;=\; j/(3\pi)
            \;\approx\; 0{,}106\,\angle\,(\pi/2),
    & c_{-3} &\;=\; -j/(3\pi)
            \;\approx\; 0{,}106\,\angle\,(-\pi/2), \\[2pt]
    c_4 &\;=\; j/(4\pi)
            \;\approx\; 0{,}080\,\angle\,(\pi/2),
    & c_{-4} &\;=\; -j/(4\pi)
            \;\approx\; 0{,}080\,\angle\,(-\pi/2), \\[2pt]
    c_5 &\;=\; j/(5\pi)
            \;\approx\; 0{,}064\,\angle\,(\pi/2),
    & c_{-5} &\;=\; -j/(5\pi)
            \;\approx\; 0{,}064\,\angle\,(-\pi/2).
\end{align*}
Los pares de índices opuestos tienen igual módulo y fases
opuestas: $|c_{-k}|=|c_k|$, $\angle c_{-k} = -\angle c_k$
---la firma de la simetría conjugada que también apareció en
el pulso, esta vez con la fase tomando un valor constante a
cada lado del cero---.

\paragraph{Espectro de líneas.}
La figura~\ref{fig:espectro_sierra} muestra el módulo $|c_k|$ y
la fase $\angle c_k$ contra el índice $k$.

\begin{figure}[H]
    \centering
    \begin{minipage}{\linewidth}
        \centering
        \includestandalone[width=0.65\linewidth]{figures/fig_espectro_sierra/fig_espectro_sierra}
        \caption{Espectro de líneas del diente de sierra con
                 $T_0 = 4$. (a) Módulo: $|c_0| = 1$ y
                 $|c_k| = 1/(\pi|k|)$ para $k\neq 0$, con
                 decaimiento $1/|k|$. (b) Fase: $0$ en $k=0$,
                 $+\pi/2$ para los $k>0$ y $-\pi/2$ para los
                 $k<0$.}
        \label{fig:espectro_sierra}
    \end{minipage}
\end{figure}

\paragraph{Función de síntesis.}
Reuniendo todo, la serie de Fourier del diente de sierra tiene
en forma compacta la expresión
\begin{equation*}
    x(t) \;=\; 1 \;+\;
        \sum_{\substack{k=-\infty \\ k\neq 0}}^{\infty}
        \frac{j}{\pi k}\,e^{j k \pi t/2},
\end{equation*}
y, sustituyendo los valores numéricos de los primeros
armónicos, se despliega como
\begin{align*}
    x(t) \;\approx\; \cdots
        & - \tfrac{j}{3\pi}\,e^{-j 3\pi t/2}
        - \tfrac{j}{2\pi}\,e^{-j\pi t}
        - \tfrac{j}{\pi}\,e^{-j\pi t/2}
        + 1 \\
        & + \tfrac{j}{\pi}\,e^{j\pi t/2}
        + \tfrac{j}{2\pi}\,e^{j\pi t}
        + \tfrac{j}{3\pi}\,e^{j 3\pi t/2}
        + \cdots
\end{align*}
La señal queda como un nivel medio $1$ más una superposición
de armónicos puramente imaginarios, antisimétrica entre
índices opuestos.

\paragraph{Comparando con el pulso.}
El contraste es directo: el pulso ---par--- da coeficientes
reales, y el diente de sierra los da puramente imaginarios
para $k\neq 0$. Ambos espectros decaen como $1/|k|$, ritmo
característico de las señales con saltos de discontinuidad.

\end{document}
